% Laget automatisk av elevark/lag_elevark.py fra sim/lysklokke.html. Endringer her blir overskrevet.
\documentclass[11pt]{article}
\usepackage{elevark}
\begin{document}
\elevarkhode{Fysikk 2 · Spesiell relativitetsteori}{Lysklokka}{https://lektorbo.no/sim/lysklokke.html}{Simuleringen}

En lysklokke er en lyspuls som går fram og tilbake mellom to speil. Hver tur ned og opp er ett tikk. Klokke A står på perrongen, og klokke B står på et tog som kjører forbi.

\section*{Oppgaver}

Jobb gjerne i par. Skriv ned gjetningen og begrunnelsen \emph{før} dere trykker Start.

\begin{oppgaveliste}
\oppgave \fase{gjett}{Gjett}Toget kjører med farten $v = 0{,}60c$. Sett fra perrongen: Går klokke B raskere, saktere eller like fort som klokke A?
\svar
\oppgave \fase{test}{Test}Kjør simuleringen til klokke A har tikket 10 ganger. Hvor mange ganger har klokke B tikket? Hvor lang tid bruker B per tikk?
\svar
\oppgave \fase{forklar}{Forklar}Slå på «Vis trekanten». Bruk Pytagoras på trekanten til å vise at tiden per tikk for B er $Δt = Δt_{0} / \sqrt{1 - v^{2}/c^{2}}$, der $Δt_{0} = 2L/c$.
\svar
\oppgave \fase{gjett}{Gjett}Hvor fort må toget kjøre for at klokke B skal gå halvparten så fort som klokke A? Regn ut svaret og test det.
\svar
\oppgave \fase{gjett}{Gjett}Bytt referansesystem til toget. Hvilken klokke går saktest nå? Er dette en selvmotsigelse?
\svar
\oppgave \fase{utfordring}{Utfordring}Et myon lever i 2,2~µs i sitt eget system. Hvor langt kommer det med farten $0{,}995c$ hvis vi ser bort fra relativitetsteorien? Hvor langt kommer det når vi tar hensyn til tidsforlengelsen?
\svar
\end{oppgaveliste}

\end{document}
